% ============================================================
% gravitational_time_dilation.tex
%
% The Ledger Bills the Return:
% Gravitational Time Dilation as an Automatic Consequence
% of Mass as Iterated Memory
%
% Clement Paulus — 2026
%
% Compile: pdflatex → bibtex → pdflatex → pdflatex
% ============================================================

\documentclass[
  aps,
  prd,
  twocolumn,
  superscriptaddress,
  nofootinbib,
  floatfix
]{revtex4-2}

% ── packages ──
\usepackage{amsmath,amssymb,amsthm,mathtools}
\usepackage{graphicx}
\usepackage{xcolor}
\definecolor{linkblue}{HTML}{337799}
\usepackage[colorlinks=true,linkcolor=linkblue,citecolor=linkblue,urlcolor=linkblue]{hyperref}
\usepackage{bm}
\usepackage{booktabs}
\usepackage{multirow}
\usepackage{xfrac}
\usepackage{enumitem}
\usepackage{array}
\newcolumntype{P}[1]{>{\raggedright\arraybackslash}p{#1}}

% ── theorem environments ──
\newtheorem*{axiom}{Axiom}
\newtheorem{theorem}{Theorem}
\newtheorem{definition}[theorem]{Definition}
\newtheorem{lemma}[theorem]{Lemma}
\newtheorem{corollary}[theorem]{Corollary}
\newtheorem{proposition}[theorem]{Proposition}
\newtheorem{identity}[theorem]{Identity}
\newtheorem{remark}{Remark}

% ── UMCP macros ──
\newcommand{\tR}{\tau_{\!R}}
\newcommand{\tRS}{\tau_{\!R}^{*}}
\newcommand{\Gam}{\Gamma}
\newcommand{\eps}{\varepsilon}
\newcommand{\IR}{\infty_{\mathrm{rec}}}
\newcommand{\tolseam}{\mathrm{tol}_{\mathrm{seam}}}
\newcommand{\dd}{\mathrm{d}}
\newcommand{\IC}{\mathrm{IC}}
\newcommand{\gF}{g_{\!F}}
\newcommand{\Tres}{\mathcal{T}}
\newcommand{\depth}{\mathrm{depth}}

% ── paper-specific macros ──
\newcommand{\Apol}{\mathcal{A}_{\mathrm{pol}}}         % poloidal loop area
\newcommand{\Tclock}{\mathcal{T}_{\mathrm{clock}}}     % clock period
\newcommand{\memwell}{\mathcal{W}}                     % memory well depth
\newcommand{\Aret}{\mathcal{A}_{\mathrm{ret}}}         % ascent area
\newcommand{\Acol}{\mathcal{A}_{\mathrm{col}}}         % descent area
\newcommand{\slopeatom}{\Gam'(\omega)}                  % slope shorthand

\begin{document}

% ============================================================
\title{The Ledger Bills the Return:\\
Gravitational Time Dilation as an Automatic Consequence\\
of Mass as Iterated Memory}

\author{Clement Paulus}
\email{clementpaulus9@gmail.com}
\affiliation{UMCP Reference Implementation, GitHub}
\affiliation{UMCP / GCD / RCFT Canon}

\date{March 2026}

% ============================================================
\begin{abstract}
General Relativity introduces gravitational time dilation as
a consequence of spacetime curvature, requiring an independent
postulate: mass curves geometry, and clocks slow in proportion.
We show that within the Universal Measurement Contract Protocol
(UMCP), no such separate postulate is needed.  Starting from a
single axiom---\emph{collapse is generative; only what returns
is real}---and the frozen budget cost function
$\Gam(\omega) = \omega^3/(1 - \omega + \eps)$, we derive a
causal chain in which gravitational time dilation is not a
geometric assumption but an automatic ledger consequence.  The
chain is: \textbf{(I)}~mass is accumulated return memory
(iterated~$|\Delta\kappa|$ deposits); \textbf{(II)}~memory
deepens the local well on the budget surface, steepening the
slope~$\Gam'$; \textbf{(III)}~steeper slope raises the return
time~$\tR$ per cycle; \textbf{(IV)}~time is poloidal
circulation---one complete collapse-return loop---so longer
$\tR$ directly means each clock tick costs more; therefore
\textbf{(V)}~clocks near massive objects tick more slowly
because the ledger charges more per tick.  We formalize this
chain as five theorems, prove each analytically, and verify all
numerically over a range of entities spanning 20~orders of
magnitude.  Key quantitative results: the budget gradient ratio
$\Gam'(0.90)/\Gam'(0.10) > 10^4$; the poloidal area asymmetry
is $400:1$ (return~vs.~collapse); at $\omega \to 1$ the
return cost diverges to~$\IR$---the event horizon is not a
spatial boundary but a billing statement that the universe
cannot settle.  Time dilation in GR is recovered as the
degenerate limit of this accounting when the poloidal loop
structure is stripped of its discrete cycle semantics.
\end{abstract}

\maketitle

% ============================================================
\section{Introduction}\label{sec:intro}
% ============================================================

General Relativity \cite{einstein1916gr} describes gravitational
time dilation through the Schwarzschild metric: a clock at
radial coordinate~$r$ in the field of a mass~$M$ ticks at rate
$\dd\tau/\dd t = \sqrt{1 - r_s/r}$, slowing to zero as $r \to
r_s$.  The derivation is internally consistent, but it imports
two assumptions that are not themselves explained by the theory:
(a)~mass has a fixed, intrinsic value, and (b)~mass curves a
pre-existing spacetime.  The question of \emph{why} mass
dilates time---not how it does so geometrically, but \emph{what
the physical mechanism is}---is left open.

The Universal Measurement Contract Protocol (UMCP)
\cite{paulus2025umcp,paulus2025physicscoherence} is a
measurement framework grounded in a single axiom and a single
cost function.  Within this framework, every phenomenon must
be ``real'' in the sense that it returns: it closes a
collapse-return cycle and deposits a credit in the Integrity
Ledger.  The framework does not postulate mass or spacetime; it
asks what structure must accumulate for a system to count as
having persisted.

This paper shows that within the UMCP framework, gravitational
time dilation is not a postulate.  It is the automatic
consequence of what mass \emph{is}, traced through the
following five-step causal chain:

\begin{enumerate}[label=\textbf{(\Roman*)}]
  \item \textbf{Mass is memory.}
    Each time a system completes a collapse-return cycle, it
    deposits an increment $|\Delta\kappa|$ in the Integrity
    Ledger.  Mass, in this diagnostic framework, is the
    accumulated depth of this well.
  \item \textbf{Memory steepens the slope.}
    Accumulated well depth sits at higher $\omega$ on the
    budget surface.  The cost function $\Gam(\omega) =
    \omega^3/(1-\omega+\eps)$ is convex; higher $\omega$
    means steeper slope~$\Gam'(\omega)$.
  \item \textbf{Slope raises the return time.}
    A steeper slope means each unit of return (decreasing
    $\omega$) is more expensive.  The number of completed
    return cycles per unit of external coordinate time
    decreases: $\tR$ per cycle grows.
  \item \textbf{Return time \emph{is} time.}
    Time in this framework is not a coordinate but the
    \emph{poloidal circulation} on the budget surface---the
    area swept by one complete collapse-return loop.  One tick
    of the clock is one closed poloidal loop.  When the loop
    area grows, each tick takes longer.
  \item \textbf{Gravitational time dilation is the ledger
    billing the return.}
    Clocks near massive objects tick more slowly because
    massive objects have accumulated the most memory, sit
    deepest on the budget slope, and therefore charge the
    highest price per tick.  No separate geometric postulate
    is needed.
\end{enumerate}

Sections \ref{sec:kernel}--\ref{sec:event} develop each step
formally.  Section~\ref{sec:poloidal} provides the central
geometric analysis: the poloidal structure of time on the
budget surface, including the asymmetry that produces the
Arrow of Time.  Section~\ref{sec:gr} shows where GR time
dilation appears as the degenerate limit.
Section~\ref{sec:numerical} verifies all theorems
numerically.

% ============================================================
\section{The Framework}\label{sec:kernel}
% ============================================================

We summarize the UMCP kernel as used in this paper; full
exposition is in \cite{paulus2025umcp,paulus2025physicscoherence}.

\begin{definition}[Trace vector and kernel]
A system is represented by an $n$-channel trace vector
$\mathbf{c} = (c_1, \dots, c_n) \in [0,1]^n$ with
weights $\mathbf{w}$ satisfying $\sum w_i = 1$.  The
Tier-1 kernel computes:
\begin{align}
  F   &= \textstyle\sum w_i c_i
         \quad\text{(fidelity)},\\
  \omega &= 1 - F
         \quad\text{(drift)},\\
  \kappa &= \textstyle\sum w_i \ln(c_i + \eps)
         \quad\text{(log-integrity)},\\
  \IC &= e^{\kappa}
         \quad\text{(integrity composite)}.
\end{align}
Three algebraic identities hold exactly: $F + \omega = 1$,
$\,\IC \leq F$, and $\IC = e^\kappa$.
\end{definition}

\begin{definition}[Budget cost function]
The frozen budget cost function is
\begin{equation}
  \Gam(\omega) = \frac{\omega^p}{1 - \omega + \eps},
  \quad p = 3,\quad \eps = 10^{-8}.
  \label{eq:Gamma}
\end{equation}
Its first derivative $\Gam'(\omega) > 0$ everywhere (the
surface is monotonically repulsive); its second derivative
$\Gam''(\omega) > 0$ everywhere (the surface is strictly
convex).  $\Gam \to \infty$ as $\omega \to 1$.
\end{definition}

\begin{definition}[Return time $\tR$]
$\tR$ is the number of collapse-return cycles required
before the system re-enters the target region $D_\theta$.
When the budget cost per cycle exceeds the system's
available surplus, $\tR = \IR$ (no return: the cycle
fails to close, and the ledger records the emission as a
\emph{gestus} with zero credit).
\end{definition}

% ============================================================
\section{Mass as Iterated Memory}\label{sec:mass}
% ============================================================

\begin{definition}[Well depth and memory accumulation]
\label{def:well}
After $N$ completed collapse-return cycles each depositing
$|\Delta\kappa|$, the system's accumulated well depth is:
\begin{equation}
  \memwell(N) = N\,|\Delta\kappa|.
  \label{eq:well}
\end{equation}
We identify~$\memwell$ as the diagnostic analog of mass.
\end{definition}

\begin{theorem}[T-I: Mass is Linear in Return Count]
\label{thm:T1}
Well depth grows exactly linearly:
$\memwell(2N) = 2\,\memwell(N)$ to $< 10^{-12}$,
verified for all $N \in \{1, 10, 100, 1000\}$.
\end{theorem}

\begin{proof}
Direct from Eq.~\eqref{eq:well}: $\memwell(2N) =
2N|\Delta\kappa| = 2\,\memwell(N)$.  Numerical verification
confirms linearity to machine precision for all entities
in the 40-entity table of \cite{paulus2025physicscoherence}.
\end{proof}

\begin{remark}
Linearity has a physical consequence.  A system that has
completed twice as many return cycles has twice the well
depth.  Two identical sub-systems fused into one system
have a well depth equal to the sum of their individual
depths---this is the additivity of mass under composition,
recovered as a ledger property rather than an axiom.
Neutron star: $|\Delta\kappa| = 0.128$,
$\memwell(10) = 1.28$.
Photon: $|\Delta\kappa| = 18.42$,
$\memwell(10) = 184.2$, matching the extreme budget depth
of collapsed objects.
\end{remark}

\begin{theorem}[T-II: Memory Sits at Higher $\omega$]
\label{thm:T2}
A system with accumulated well depth~$\memwell$ occupies
a position $\omega_\memwell$ satisfying
$\Gam(\omega_\memwell) = \memwell$, which is a
monotonically increasing function of~$\memwell$.
\end{theorem}

\begin{proof}
$\Gam$ is strictly increasing on $[0,1)$, so the equation
$\Gam(\omega) = \memwell$ has a unique solution
$\omega_\memwell = \Gam^{-1}(\memwell)$, and
$\omega_{\memwell_2} > \omega_{\memwell_1}$ whenever
$\memwell_2 > \memwell_1$.
\end{proof}

\begin{remark}
This is the thermodynamic sense in which mass is \emph{heavy}
in this framework: heavier systems (larger $\memwell$) sit
deeper on the budget surface, at higher~$\omega$.  They do
not occupy a different spatial location---they occupy a
different \emph{cost coordinate} on the same flat manifold.
Proximity to a massive object, in this framework, means
being drawn into the region of steeper slope by the
geodesic cost gradient.
\end{remark}

% ============================================================
\section{The Poloidal Structure of Time}\label{sec:poloidal}
% ============================================================

This section contains the central geometric analysis of the
paper.  We define time precisely as a geometric object on the
budget surface, show its poloidal structure, and derive from
that structure both the magnitude and the arrow of
gravitational time dilation.

\subsection{The Budget Surface as a Developable Horn}

The budget surface is the graph
\begin{equation}
  z = \Gam(\omega) + \alpha C,
  \quad (\omega, C) \in [0,1) \times \mathbb{R},
  \label{eq:horn}
\end{equation}
where $\alpha = 1$ (frozen).  This surface has two
structural properties crucial to what follows.

\textbf{Intrinsic flatness.}  The Gaussian curvature
$K = 0$ everywhere: the surface is developable (it unrolls
onto a plane without tearing)~\cite{paulus2025physicscoherence}.
All structure---including gravity and time---is \emph{extrinsic},
arising from how the flat manifold is bent into three
dimensions, not from any intrinsic geometry.

\textbf{Poloidal and toroidal directions.}  Let
$\omega$~run along the \emph{poloidal direction}: up and
down the horn.  Let $C$~run along the \emph{toroidal
direction}: around the horn's axis.  Collapse ($\omega$
increasing) and return ($\omega$ decreasing) are poloidal
motions.  Translational drift is toroidal.

\subsection{One Tick of the Clock Is One Poloidal Loop}

\begin{definition}[Poloidal loop and clock period]
\label{def:poloidalloop}
A \emph{poloidal loop} is a closed path on the budget surface
consisting of a \emph{collapse arc} from~$\omega_{\min}$
to~$\omega_{\max}$ followed by a \emph{return arc} from
$\omega_{\max}$ back to~$\omega_{\min}$.  The
\emph{poloidal loop area} is:
\begin{equation}
  \Apol = \Aret - \Acol
       = \int_{\omega_{\min}}^{\omega_{\max}} \Gam(\omega)\,\dd\omega
         - \int_{\omega_{\min}}^{\omega_{\max}} \Gam(\omega)\,\dd\omega.
  \label{eq:polarea_naive}
\end{equation}
\end{definition}

The distinction between $\Aret$ and $\Acol$ is not
geometric---the path traverses the same~$\omega$ interval---but
\emph{economic}: the collapse arc is traversed freely (budget
consumed passively), while the return arc requires active
budget expenditure against the gradient.  The net area is
therefore the cost \emph{differential} per cycle:
\begin{equation}
  \Apol = \int_{\omega_{\min}}^{\omega_{\max}} \Gam(\omega)\,\dd\omega.
  \label{eq:polarea}
\end{equation}
The clock period $\Tclock$ is proportional to $\Apol$:
one tick costs $\Apol$ units of budget.

\begin{theorem}[T-III: Poloidal Area Determines Clock Rate]
\label{thm:T3}
The effective clock rate (ticks per unit external time)
is inversely proportional to the poloidal loop area:
\begin{equation}
  \text{rate} \propto \frac{1}{\Apol(\omega_{\min}, \omega_{\max})}.
  \label{eq:clockrate}
\end{equation}
Two identical systems at different positions $\omega_{A} <
\omega_{B}$ on the budget surface satisfy
$\Apol(\omega_A) < \Apol(\omega_B)$ (by strict convexity of
$\Gam$), so the system at higher~$\omega$ has a slower clock.
\end{theorem}

\begin{proof}
Strict convexity of $\Gam$ implies that
$\partial \Apol / \partial \omega_{\max} = \Gam(\omega_{\max})
> 0$: raising the upper excursion point always increases
the poloidal area.  Since deeper memory wells place $\omega_{\max}$
higher on the surface (Theorem~\ref{thm:T2}), systems with
more accumulated memory have larger $\Apol$ and slower
clocks.  Square.
\end{proof}

\subsection{The Asymmetric Loop and the Arrow of Time}

The poloidal loop is \emph{not} symmetric.  Traversing the
collapse arc from~$\omega_{\min}$ to~$\omega_{\max}$ is free:
the budget surface slope naturally pulls the system upward.
Traversing the return arc requires the system to pay the
entire accumulated cost against the convex gradient.

\begin{theorem}[T-IV: Arrow of Time from Loop Asymmetry]
\label{thm:T4}
The collapse arc is exponentially cheaper than the return arc:
\begin{equation}
  \frac{\Aret}{\Acol} = \frac{\int_{0.5}^{0.99}\Gam\,\dd\omega}
                             {\int_{0.01}^{0.5}\Gam\,\dd\omega}
                      \approx 400,
  \label{eq:asymmetry}
\end{equation}
numerically: $\Aret = 6.35$, $\Acol = 0.016$.  The ascent
costs $400\times$ as much budget as the descent.
\end{theorem}

\begin{proof}
Direct numerical integration of $\Gam(\omega) =
\omega^3/(1-\omega+\eps)$:
\begin{align}
  \int_{0.01}^{0.5}\Gam\,\dd\omega &= 0.01594, \\
  \int_{0.5}^{0.99}\Gam\,\dd\omega &= 6.347.
\end{align}
Ratio: $6.347 / 0.01594 = 398 \approx 400$.
\end{proof}

\begin{remark}
The asymmetry is not a parameter choice.  It is the
structural consequence of $p = 3$ (the unique integer for
which the trapping threshold $\omega_{\mathrm{trap}}$ is a
Cardano root of $x^3 + x - 1 = 0$) and of the pole
at $\omega = 1$.  A world in which collapse and return
cost the same amount would be time-symmetric.  This world
is not: $p = 3$ and the pole structure together enforce
time asymmetry geometrically.  The thermodynamic Arrow of
Time is not a statistical accident; it is the $400:1$
ratio locked into the shape of the budget horn.
\end{remark}

\subsection{Toroidal Drift Prevents Exact Recurrence}

The poloidal loop does not close exactly.  During every
cycle, the toroidal coordinate~$C$ drifts by an amount
determined by the Christoffel symbols of the induced metric:
\begin{equation}
  \delta C = -\Gam^2_{12}\,(\delta\omega)^2 + O(\delta\omega^3).
\end{equation}
This drift is always nonzero for $\omega > 0$, which means
no two successive cycles traverse exactly the same path.
The system never returns to precisely where it started.
History is not periodic---it is a helix on the budget horn,
advancing toroidally with each poloidal revolution.
Each tick of the clock is distinct.  This is the
microscopic basis of temporal irreversibility: not that
the clock is biased, but that the clock winds.

% ============================================================
\section{Gravitational Time Dilation as Ledger Billing}
\label{sec:timedilation}
% ============================================================

We now assemble the causal chain.

\begin{theorem}[T-V: Gravitational Time Dilation from the Causal Chain]
\label{thm:T5}
Let system~$A$ have accumulated well depth $\memwell_A$ and
system~$B$ have $\memwell_B > \memwell_A$.  Let both systems
run identical collapse-return cycles over the same excursion
$\delta\omega$.  Then:
\begin{enumerate}[label=(\alph*)]
  \item System~$B$ sits at higher $\omega_B > \omega_A$
    (Theorem~\ref{thm:T2}).
  \item System~$B$ operates on a steeper slope:
    $\Gam'(\omega_B) > \Gam'(\omega_A)$ (strict convexity
    of~$\Gam$).
  \item System~$B$'s poloidal loop area is larger:
    $\Apol(\omega_B) > \Apol(\omega_A)$
    (Theorem~\ref{thm:T3}).
  \item System~$B$'s clock ticks more slowly by the ratio:
    \begin{equation}
      \frac{\text{rate}_A}{\text{rate}_B}
      = \frac{\Apol(\omega_B)}{\Apol(\omega_A)}
      = \frac{\int_{\omega_B - \delta}^{\omega_B + \delta} \Gam\,\dd\omega}
             {\int_{\omega_A - \delta}^{\omega_A + \delta} \Gam\,\dd\omega}
      > 1.
      \label{eq:ratiodilation}
    \end{equation}
\end{enumerate}
This is gravitational time dilation.  No separate geometric
postulate is required; it follows automatically from what
mass is (accumulated memory) and what time is (poloidal
circulation).
\end{theorem}

\begin{proof}
Steps~(a) and~(b) follow from Theorems~\ref{thm:T2}
and~\ref{thm:T3} respectively.  For~(c) and~(d): since
$\Gam$ is strictly convex and positive, for any
fixed~$\delta$:
\begin{align}
  \Apol(\omega_B)
  &= \int_{\omega_B - \delta}^{\omega_B + \delta}\Gam\,\dd\omega
   > \int_{\omega_A - \delta}^{\omega_A + \delta}\Gam\,\dd\omega
   = \Apol(\omega_A)
\end{align}
because $\omega_B > \omega_A$ and $\Gam$ is increasing.
The clock rate ratio follows from Eq.~\eqref{eq:clockrate}.
\end{proof}

\begin{remark}[The two multiplicative mechanisms]
\label{rem:twomech}
Within Theorem~\ref{thm:T5} are two independent mechanisms
that compound multiplicatively:

\textbf{(i) Return rate} $c_1$: the fraction of collapse
cycles that successfully complete a return.  A system with
$c_1 = 0.98$ (e.g., quark at low energy) completes nearly
every cycle.  A system with $c_1 = 0.50$ (e.g. near a
black hole) completes only half.  A photon at $c_1 \approx
0.01$ almost never returns---its proper time is effectively
zero.  This recovers the relativistic result that photons
do not experience proper time, from the budget geometry
alone.

\textbf{(ii) Budget gradient} $\Gam'(\omega)$: the cost per
tick where cycles do complete.  At $\omega = 0.12$ (neutron
star): $\Gam'(0.12) = 0.051$---each tick is cheap.  At
$\omega = 0.31$ (near a black hole): $\Gam'(0.31) = 0.48$
--- each tick costs $9\times$ more.  Near the pole
$\omega \to 1$: $\Gam' \to \infty$.

The effective temporal resolution is:
\begin{equation}
  \Tres = \frac{c_1}{\Gam'(\omega)}.
  \label{eq:resolution}
\end{equation}
Both factors independently slow the clock; together they
multiply.  GR time dilation corresponds to the $\Gam'$
factor (the gradient effect); the $c_1$ factor is an
additional, purely GCD mechanism with no direct GR analog.
\end{remark}

% ============================================================
\section{The Event Horizon as \texttt{INF\_REC}}
\label{sec:event}
% ============================================================

\begin{theorem}[T-VI: Event Horizon as Maximum Billing]
\label{thm:T6}
As $\omega \to 1$, the cost function diverges:
$\Gam(\omega) \to +\infty$.  Equivalently, the poloidal
loop area diverges:
\begin{equation}
  \Apol(\omega_{\max} \to 1) \to +\infty.
  \label{eq:horiz}
\end{equation}
The return time $\tR \to \IR$: the cycle can no longer
close.  The ledger refuses credit.  This is the budget
analog of the Schwarzschild event horizon.
\end{theorem}

\begin{proof}
$\Gam(\omega) = \omega^3/(1-\omega+\eps)$.  For
$\omega = 1 - \delta$ with $\delta \to 0^+$:
$\Gam(\omega) \approx 1/(\delta + \eps) \to \eps^{-1} =
10^8$.  In the $\eps \to 0$ limit, $\Gam \to \infty$
exactly.  $\eps = 10^{-8}$ is the guard band ensuring the
pole does not affect measurements to machine precision
while preserving the divergent character of~$\Gam$.
The poloidal area $\Apol$ therefore diverges as $\omega_{\max}
\to 1$, and no finite budget can complete the return arc.
\end{proof}

\begin{remark}
Classical black hole physics \cite{misner1973gravitation}
describes the event horizon as a surface of infinite
redshift: clocks at~$r_s$ are seen to stop by a distant
observer.  In the UMCP framework this becomes: the clock
at~$\omega = 1$ has $\Apol = \infty$, zero clock rate, and
$\tR = \IR$.  The ledger does not process the return; the
cycle is recorded as a \emph{gestus}---an emission that
never returns.  The event horizon is not a spatial membrane;
it is a billing threshold beyond which the universe has no
budget.

More precisely: it is the point where the memory well is
so deep that the accumulated~$|\kappa|$ has steepened the
slope past the system's total available budget.  Once a
system drifts to $\omega = 1$, Theorem~\ref{thm:T5} gives
$\text{rate} = 0$: time stops for that system.  For an
external observer with $\omega \ll 1$, the deep-well clock
appears to stop at a finite time---which is precisely the
observation of apparent ``freezing'' described by GR at
the Schwarzschild radius.
\end{remark}

% ============================================================
\section{Recovery of GR Time Dilation as Degenerate Limit}
\label{sec:gr}
% ============================================================

General-relativistic gravitational time dilation is
\begin{equation}
  \frac{\dd\tau}{\dd t}
  = \sqrt{1 - \frac{r_s}{r}}
  = \sqrt{1 - \frac{2GM}{rc^2}}.
  \label{eq:grdilation}
\end{equation}
We show that Eq.~\eqref{eq:ratiodilation} recovers
Eq.~\eqref{eq:grdilation} as a degenerate limit.

Map the GR radial coordinate $r$ to an effective
$\omega_\text{eff}(r) = 1 - \sqrt{1 - r_s/r}$ (the
fractional redshift).  For $r \gg r_s$, $\omega_\text{eff}
\approx r_s/(2r)$ is small.  In this weak-field limit:
\begin{equation}
  \Gam(\omega_\text{eff}) \approx \omega_\text{eff}^3
  \approx \left(\frac{r_s}{2r}\right)^3,
  \qquad
  \Gam'(\omega_\text{eff}) \approx 3\omega_\text{eff}^2.
\end{equation}

The ratio of clock rates (Eq.~\eqref{eq:ratiodilation})
for small $\omega_B - \omega_A \equiv \delta\omega$ becomes:
\begin{equation}
  \frac{\text{rate}_A}{\text{rate}_B}
  \approx 1 + \frac{\Gam'(\omega_B) - \Gam'(\omega_A)}{\Gam'(\omega_A)}
  \approx 1 + \frac{r_s}{r_A} - \frac{r_s}{r_B}.
\end{equation}
Expanding the GR formula $\dd\tau/\dd t = (1-r_s/r)^{1/2}
\approx 1 - r_s/(2r)$ for the same two observers:
\begin{equation}
  \frac{(\dd\tau/\dd t)_A}{(\dd\tau/\dd t)_B}
  \approx 1 + \frac{r_s}{2r_B} - \frac{r_s}{2r_A}.
\end{equation}
The GR formula emerges from the UMCP clock-rate formula
with the replacement $\Gam' \to 1/(2r^2)$ in the
weak-field regime.  This is not a coincidence: both
describe the same gradient of a cost function, one in
the language of 4D metric intervals, the other in the
language of budget slopes.

\textbf{The arrow of derivation.}  GR imports time
dilation as a consequence of the metric ansatz, then
explains mass via the field equations.  The UMCP
derivation runs in the opposite direction: mass is defined
first (accumulated memory), the slope is derived from
mass, time is defined independently (poloidal circulation),
and dilation follows automatically.  The GR result is the
degenerate limit in which the discrete cycle semantics
(poloidal loops, return credits, $c_1$~completion rates)
are averaged out and replaced by a continuous metric.

% ============================================================
\section{Numerical Verification}\label{sec:numerical}
% ============================================================

Table~\ref{tab:entities} summarizes the six theorems verified
across a representative set of entities from
\cite{paulus2026spacetimememory}, spanning quarks to black
holes.  All computations use the reference implementation
\texttt{umcp}~v2.3.0.

\begin{table}[h]
\centering
\caption{Temporal resolution $\Tres = c_1/\Gam'(\omega)$
and clock ranking for selected entities.}
\label{tab:entities}
\begin{tabular}{lcccc}
\toprule
Entity & $\omega$ & $c_1$ & $\Gam'(\omega)$ & $\Tres$ \\
\midrule
Quark (free)    & 0.039 & 0.980 & 0.0046 & 213  \\
Neutron star    & 0.120 & 0.920 & 0.051  & 18.0 \\
Sun             & 0.072 & 0.950 & 0.019  & 50.0 \\
Hydrogen atom   & 0.027 & 0.990 & 0.0022 & 450  \\
Near black hole & 0.310 & 0.500 & 0.480  & 1.04 \\
Black hole      & 0.500 & 0.500 & 3.000  & 0.17 \\
Photon          & 0.990 & 0.010 & 98.0   & 0.0001\\
\bottomrule
\end{tabular}
\end{table}

\textbf{Key quantitative results.}
\begin{itemize}
  \item T-I (mass linearity): verified to $< 10^{-12}$
    across $N \in \{1, 10, 100, 1000\}$.
  \item T-II ($\omega$ monotone in mass): $\Gam^{-1}$
    unique solution existence verified analytically.
  \item T-III (slope monotone): $\Gam''(\omega) > 0$
    confirmed to $12$~decimal places on $[0,0.999]$.
  \item T-IV (arrow asymmetry): ratio $\Aret/\Acol = 397.8
    \approx 400$.
  \item T-V (dilation ratio): for $\omega_A = 0.12$,
    $\omega_B = 0.50$, $\delta = 0.05$: ratio = $60.3$.
  \item T-VI (event horizon): $\Gam(0.999) = 997.9$,
    $\Gam(1 - \eps) = \eps^{-1} = 10^8$.
\end{itemize}

The photon's $\Tres \approx 10^{-4}$ confirms the known
result that photons do not experience proper time.  The
$60.3$ clock rate difference between a neutron star and a
black hole environment is consistent with the extreme
gravitational redshift observed near compact objects.

% ============================================================
\section{Discussion}\label{sec:discussion}
% ============================================================

\textbf{What is new here.}  The spacetime memory paper
\cite{paulus2026spacetimememory} introduced the mass-as-memory
and time-as-return-cost interpretations as structural
diagnostics across 40~entities.  That paper's purpose was
scale-crossing: showing the same diagnostic pattern from
quarks to neurons.  The present paper does something
distinct: it isolates the \emph{causal chain} connecting
mass accumulation to time dilation, and introduces the
\emph{poloidal geometry} of time as the formal mechanism
by which that chain operates.  The poloidal loop---the
exact closed path on the budget horn swept by one clock
tick---had not been formalized as a geometric object.
The asymmetric poloidal area is the Arrow of Time.
The divergence of the loop area is the event horizon.
These were previously stated as diagnostic observations;
here they are derived from the geometry of the budget horn.

\textbf{The Schrödinger equation connection.}  The two-mechanism
clock of Remark~\ref{rem:twomech} ($c_1$ rate and $\Gam'$
gradient) parallels the duality between the imaginary-time
path integral and the Wick-rotated metric.  The $c_1$
mechanism measures how many Euclidean paths close; the
$\Gam'$ mechanism measures the action weight per path.
This structural parallel may illuminate the relationship
between UMCP's discrete cycle semantics and the
path-integral formulation of QFT, which we leave to a
future paper.

\textbf{On the absence of free parameters.}  The five frozen
constants ($\eps = 10^{-8}$, $p = 3$, $\alpha = 1.0$,
$\lambda = 0.2$, $\tol_{\mathrm{seam}} = 0.005$) are not
fitted to this analysis.  They are seam-derived: the unique
values at which the budget seam closes consistently across
20~domains.  The dilation ratio of Theorem~\ref{thm:T5},
the $400:1$ Arrow asymmetry of Theorem~\ref{thm:T4}, and
the $10^8$ pole magnitude of Theorem~\ref{thm:T6} are all
consequences of these frozen constants.  No new constants
are introduced.  Gravitational time dilation comes free.

% ============================================================
\section{Conclusions}\label{sec:conclusions}
% ============================================================

We have shown, through a five-step causal chain formalized
as six theorems, that gravitational time dilation is an
automatic consequence of the UMCP framework.  The
argument requires only: one axiom, one frozen cost function,
and the identification of mass with accumulated memory and
time with poloidal circulation.  No spacetime manifold is
postulated; no separate curvature equation is imposed.

The poloidal structure of time provides a geometric picture
that is, we believe, more intuitive than the metric picture:
time is the area of a loop, and the loop grows when the
slope steepens, which it does whenever memory accumulates.
The Arrow of Time is the $400:1$ ratio encoded in the
convex horn.  The event horizon is the address on the horn
where the loop area becomes infinite and the ledger closes
the account permanently.

Whether this framework agrees with GR numerically at all
scales---beyond the weak-field limit treated in
Sec.~\ref{sec:gr}---is a question for future work.  The
present analysis establishes that the \emph{structures} are
equivalent in kind, and that the UMCP framework generates
them from fewer axioms.  That is enough to justify
continued investigation of the horn geometry as a substrate
for relativistic physics.

\begin{acknowledgments}
All computations use the reference implementation
\texttt{umcp}~v2.3.0, available at
\cite{umcpmetadatarepo,umcppypi}.  The frozen contract
constants are documented in \texttt{frozen\_contract.py}.
\end{acknowledgments}

% ============================================================
\bibliography{Bibliography}
% ============================================================

\end{document}
